The source geometry and setup parameters used for observation are:
\[
L = 0.36\,\text{cm}, \quad r_o = 1.5\,\text{cm}, \quad \theta_o = 90^\circ
\]

\begin{figure}[H]
	\centering
	\includegraphics[width=\textwidth]{\detokenize{Figures/Experiment/Observation table.pdf}}
\end{figure}
\subsection{\texorpdfstring{$gL(r)$ dependency on $r$}{gL(r) dependency on r}}
The variation of the geometry function \( g_L(r) \) with respect to the distance \( r \) from the source is depicted in the following graph:
\begin{figure}[H]
	\centering
	\includegraphics[width=0.7\textwidth]{\detokenize{Figures/Experiment/gL(r) vs r.pdf}}
\end{figure}

\subsection{\texorpdfstring{$F(r, \theta)$ on different angles($\theta$)}{F(r, theta) on different angles(theta)}}
The variation of the anisotropy function \( F(r, \theta) \) with respect to different angles \( \theta \) is shown in the following graph:
\begin{figure}[H]
\centering
\begin{minipage}[t]{0.32\textwidth}
	\centering
	\includegraphics[width=\textwidth]{\detokenize{Figures/Experiment/1.5 cm.pdf}}\\[-1mm]
	{\small \(r = 1.5\,\text{cm}\)}
\end{minipage}\hfill
\begin{minipage}[t]{0.32\textwidth}
	\centering
	\includegraphics[width=\textwidth]{\detokenize{Figures/Experiment/3 cm.pdf}}\\[-1mm]
	{\small \(r = 3\,\text{cm}\)}
\end{minipage}\hfill
\begin{minipage}[t]{0.32\textwidth}
	\centering
	\includegraphics[width=\textwidth]{\detokenize{Figures/Experiment/4.5 cm.pdf}}\\[-1mm]
	{\small \(r = 4.5\,\text{cm}\)}
\end{minipage}
\caption{Anisotropy function \(F(r, \theta)\) at different radial distances}
\end{figure}
